%! The Victory of Orthogonality — Unit 7, Lesson 7.4 — Warm-Up
\documentclass[10pt]{article}
\usepackage{linalg-article}
\usepackage{linalg-boxes}
\newcommand{\TallMath}[1]{$\displaystyle #1\rule[-1.4em]{0pt}{3.2em}$}
\newcommand{\uu}{\mathbf{u}}
\newcommand{\vv}{\mathbf{v}}
\newcommand{\xx}{\mathbf{x}}
\newcommand{\qq}{\mathbf{q}}
\newcommand{\T}{^{\top}}

\begin{document}

\pageheader{Unit 7, Lesson 7.4}{Warm-Up}
\namedateperiod

\vspace{0.12in}

\begin{notesbox}{Getting Ready: the Set-Up Moves of Today's Idea}
\small These three quick checks --- testing perpendicular unit columns, forming $Q\T Q$, and measuring length
after $Q$ --- \emph{are} the opening moves of today's lesson. We use the matrix that has appeared all unit:
\[
  Q=\tfrac1{\sqrt2}\begin{bmatrix} 1 & 1\\ 1 & -1 \end{bmatrix},\qquad
  \text{columns }\ \qq_1=\tfrac1{\sqrt2}\begin{bmatrix} 1\\1 \end{bmatrix},\ \ \qq_2=\tfrac1{\sqrt2}\begin{bmatrix} 1\\-1 \end{bmatrix}.
\]

\vspace{2pt}
\textbf{1. Perpendicular unit columns? (Unit 4).} Use dot products. Each column's length is
$\sqrt{\qq\cdot\qq}$, and two vectors are perpendicular when their dot product is $0$.
\[
  \qq_1\cdot\qq_1=\blank{1.0cm},\qquad \qq_2\cdot\qq_2=\blank{1.0cm},\qquad \qq_1\cdot\qq_2=\blank{1.0cm}.
\]

\textbf{2. Form $Q\T Q$ (Unit 1 \& Unit 6).} Each entry of $Q\T Q$ is a dot product of two columns of $Q$:
\[
  Q\T Q=\begin{bmatrix}\rule{0pt}{1.1em}\blank{0.9cm} & \blank{0.9cm}\\[2pt]\blank{0.9cm} & \blank{0.9cm}\end{bmatrix}.
  \qquad\text{So $Q^{-1}=$ }\blank{1.6cm}.
\]

\textbf{3. Length after $Q$ (Unit 4).} Let $\xx=(3,1)$, so $|\xx|^2=10$. Compute $Q\xx$ and its length squared.
\[
  Q\xx=\tfrac1{\sqrt2}\begin{bmatrix}\rule{0pt}{1.1em}\blank{0.7cm}\\[2pt]\blank{0.7cm}\end{bmatrix},\qquad
  |Q\xx|^2=\blank{1.2cm}.
\]
Did $Q$ change the length of $\xx$?\ \ \blank{2.4cm}
\end{notesbox}

\end{document}
